
\chapter{2016年全国I卷（文）}

\section{单选题}

\begin{problem}
设集合 \(A=\{1,3,5,7\}\)，\(B=\{x\mid 2\le x\le 5\}\)，则 \(A\cap B=\)（\quad）

\choices
    {\(\{1,3\}\)}
    {\(\{3,5\}\)}
    {\(\{5,7\}\)}
    {\(\{1,7\}\)}
\end{problem}

\begin{answer}
B
\end{answer}

\begin{solution}
集合 \(B\) 中包含区间 \([2,5]\) 内的所有实数，因此 \(A\) 中属于 \(B\) 的元素是 \(3\) 和 \(5\)，从而 \(A\cap B=\{3,5\}\)，故选 B.
\end{solution}

\begin{problem}
设 \((1+2\i)(a+\i)\) 的实部与虚部相等，其中 \(a\) 为实数，则 \(a=\)（\quad）

\choices
    {\(-3\)}
    {\(-2\)}
    {\(2\)}
    {\(3\)}
\end{problem}

\begin{answer}
A
\end{answer}

\begin{solution}
展开复数得 \((1+2\i)(a+\i)=a+\i+2a\i+2\i^2=(a-2)+(2a+1)\i\). 实部与虚部相等，所以 \(a-2=2a+1\)，解得 \(a=-3\)，故选 A.
\end{solution}

\begin{problem}
为美化环境，从红，黄，白，紫 4 种颜色的花中任选 2 种花种在一个花坛中，余下的 2 种花种在另一个花坛中，则红色和紫色的花不在同一花坛的概率是（\quad）

\choices
    {\(\dfrac13\)}
    {\(\dfrac12\)}
    {\(\dfrac23\)}
    {\(\dfrac56\)}
\end{problem}

\begin{answer}
C
\end{answer}

\begin{solution}
从 4 种花中任选 2 种放入指定的一个花坛，基本事件总数为 \(\mathrm C_4^2=6\). 若红色和紫色不在同一花坛，所选的 2 种花中必须恰有一种是红色或紫色：先选红、紫中的一种有 \(2\) 种，再从黄、白中选一种有 \(2\) 种，所以有 \(4\) 个有利事件. 因此所求概率为 \(\dfrac{4}{6}=\dfrac23\)，故选 C.
\end{solution}

\begin{problem}
\(\triangle ABC\) 的内角 \(A\)，\(B\)，\(C\) 的对边分别为 \(a\)，\(b\)，\(c\). 已知 \(a=\sqrt5\)，\(c=2\)，\(\cos A=\dfrac23\)，则 \(b=\)（\quad）

\choices
    {\(\sqrt2\)}
    {\(\sqrt3\)}
    {\(2\)}
    {\(3\)}
\end{problem}

\begin{answer}
D
\end{answer}

\begin{solution}
由余弦定理，\(a^2=b^2+c^2-2bc\cos A\). 代入已知数据得 \(5=b^2+4-2\cdot b\cdot2\cdot\dfrac23\)，即 \(3b^2-8b-3=0\). 解得 \(b=3\) 或 \(b=-\dfrac13\)，但边长必须为正，所以 \(b=3\)，故选 D.
\end{solution}

\begin{problem}
直线 \(l\) 经过椭圆的一个顶点和一个焦点，若椭圆中心到 \(l\) 的距离为其短轴长的 \(\dfrac14\)，则该椭圆的离心率为（\quad）

\choices
    {\(\dfrac13\)}
    {\(\dfrac12\)}
    {\(\dfrac23\)}
    {\(\dfrac34\)}
\end{problem}

\begin{answer}
B
\end{answer}

\begin{solution}
设椭圆的半长轴、半短轴和焦距分别为 \(a_0\)，\(b_0\)，\(c_0\)，则 \(a_0^2=b_0^2+c_0^2\). 若所取顶点是长轴顶点，连接该顶点与焦点得到的就是长轴，直线到中心的距离为零，与题意不符，所以所取顶点是短轴顶点. 取短轴顶点为 \(V=(0,b_0)\)，焦点为 \(F=(c_0,0)\)，则直线 \(VF\) 的方程为 \(b_0x+c_0y-b_0c_0=0\).

椭圆中心到该直线的距离为 \(\dfrac{b_0c_0}{\sqrt{b_0^2+c_0^2}}=\dfrac{b_0c_0}{a_0}\). 又因短轴长为 \(2b_0\)，由题意有 \(\dfrac{b_0c_0}{a_0}=\dfrac{b_0}{2}\)，所以 \(\dfrac{c_0}{a_0}=\dfrac12\). 因此离心率为 \(\dfrac12\)，故选 B.
\end{solution}

\begin{problem}
将函数 \(y=2\sin\!\left(2x+\dfrac{\pi}{6}\right)\) 的图象向右平移 \(\dfrac14\) 个周期后，所得图象对应的函数为（\quad）

\choices
    {\(y=2\sin\!\left(2x+\dfrac{\pi}{4}\right)\)}
    {\(y=2\sin\!\left(2x+\dfrac{\pi}{3}\right)\)}
    {\(y=2\sin\!\left(2x-\dfrac{\pi}{4}\right)\)}
    {\(y=2\sin\!\left(2x-\dfrac{\pi}{3}\right)\)}
\end{problem}

\begin{answer}
D
\end{answer}

\begin{solution}
原函数的周期为 \(T=\dfrac{2\pi}{2}=\pi\). 向右平移 \(\dfrac14T=\dfrac\pi4\) 后，将原式中的 \(x\) 换成 \(x-\dfrac\pi4\)，得到 \(y=2\sin\!\left(2\left(x-\dfrac\pi4\right)+\dfrac\pi6\right)=2\sin\!\left(2x-\dfrac\pi3\right)\)，故选 D.
\end{solution}

\begin{problem}
如图，某几何体的三视图是三个半径相等的圆及每个圆中两条相互垂直的半径. 若该几何体的体积是 \(28\pi/3\)，则它的表面积是（\quad）

\FigureLayoutDeclare{img/2016/national_paper_1_liberal/q7_fig1.png}{4.0cm}{34bp 41bp 272bp 0bp}
\bitmapfigure[max width=6.0cm]{img/2016/national_paper_1_liberal/q7_fig1.png}

\choices
    {\(17\pi\)}
    {\(18\pi\)}
    {\(20\pi\)}
    {\(28\pi\)}
\end{problem}

\begin{answer}
A
\end{answer}

\begin{solution}
由三视图可知，该几何体是半径为 \(r\) 的球去掉一个八分之一球体，因而其体积为球体积的 \(\dfrac78\). 所以 \(\dfrac78\cdot\dfrac43\pi r^3=\dfrac{28\pi}{3}\)，解得 \(r=2\).

它的曲面部分是球面的 \(\dfrac78\)，面积为 \(\dfrac78\cdot4\pi r^2=14\pi\)；去掉八分之一球体后新增三个互相垂直的四分之一圆面，每个面积为 \(\dfrac14\pi r^2=\pi\). 所以表面积为 \(14\pi+3\pi=17\pi\)，故选 A.
\end{solution}

\begin{problem}
若 \(a>b>0\)，\(0<c<1\)，则（\quad）

\choices
    {\(\log_a c<\log_b c\)}
    {\(\log_c a<\log_c b\)}
    {\(a^c<b^c\)}
    {\(c^a>c^b\)}
\end{problem}

\begin{answer}
B
\end{answer}

\begin{solution}
因为 \(0<c<1\)，函数 \(y=\log_c x\) 在正数范围内单调递减，所以由 \(a>b>0\) 得 \(\log_c a<\log_c b\). 同时，指数函数 \(y=x^c\) 在正数范围内递增，故 \(a^c>b^c\)；函数 \(y=c^x\) 递减，故 \(c^a<c^b\). 因此只有 B 正确，故选 B.
\end{solution}

\begin{problem}
函数 \(y=2x^2-\e^{|x|}\) 在 \([-2,2]\) 的图象大致为（\quad）

\choices
    {\choicebitmap[width=0.15\paperwidth]{img/2016/national_paper_1_liberal/q9_fig3.png}}
    {\choicebitmap[width=0.15\paperwidth]{img/2016/national_paper_1_liberal/q9_fig1.png}}
    {\choicebitmap[width=0.15\paperwidth]{img/2016/national_paper_1_liberal/q9_fig4.png}}
    {\choicebitmap[width=0.15\paperwidth]{img/2016/national_paper_1_liberal/q9_fig2.png}}
\end{problem}

\begin{answer}
D
\end{answer}

\begin{solution}
函数为偶函数，且 \(f(0)=-1\)，\(f(1)=2-\e<0\)，\(f(2)=8-\e^2>0\)，所以图象关于 \(y\) 轴对称，原点附近在 \(x\) 轴下方，且在两端上方. 对 \(x\ge0\)，有 \(f'(x)=4x-\e^x\). 令 \(h(x)=4x-\e^x\)，则 \(h'(x)=4-\e^x\)，所以 \(h\) 在 \([0,\ln4]\) 上递增，在 \([\ln4,2]\) 上递减. 又 \(h(0)=-1<0\)，\(h(1)=4-\e>0\)，\(h(2)=8-\e^2>0\)，故 \(h\) 在 \([0,2]\) 内恰有一个零点，且该零点在 \((0,1)\) 内. 因此 \(f\) 在正半轴先减后增，在正半轴有一个极小值，并在 \((1,2)\) 内穿过 \(x\) 轴，负半轴部分由偶性对称得到. 符合上述特征的图象为 D，故选 D.
\end{solution}

\begin{problem}
执行下面的程序框图，如果输入的 \(x=0\)，\(y=1\)，\(n=1\)，则输出 \(x\)，\(y\) 的值满足（\quad）

\bitmapfigure[max width=6.0cm]{img/2016/national_paper_1_liberal/q10_fig1.png}

\choices
    {\(y=2x\)}
    {\(y=3x\)}
    {\(y=4x\)}
    {\(y=5x\)}
\end{problem}

\begin{answer}
C
\end{answer}

\begin{solution}
第一次循环时 \(n=1\)，更新后 \(x=0+\dfrac{1-1}{2}=0\)，\(y=1\cdot1=1\)，此时 \(x^2+y^2=1<36\)，令 \(n=2\). 第二次循环后 \(x=0+\dfrac{2-1}{2}=\dfrac12\)，\(y=2\)，此时 \(x^2+y^2=\dfrac14+4<36\)，令 \(n=3\). 第三次循环后 \(x=\dfrac12+\dfrac{3-1}{2}=\dfrac32\)，\(y=3\cdot2=6\)，且 \(x^2+y^2=\dfrac94+36>36\)，循环结束. 于是 \(y=6=4\cdot\dfrac32=4x\)，故选 C.
\end{solution}

\begin{problem}
平面 \(\alpha\) 过正方体 \(ABCD-A_1B_1C_1D_1\) 的顶点 \(A\)，\(\alpha\parallel\) 平面 \(CB_1D_1\)，\(\alpha\cap\) 平面 \(ABCD=m\)，\(\alpha\cap\) 平面 \(ABB_1A_1=n\)，则 \(m\)，\(n\) 所成角的正弦值为（\quad）

\choices
    {\(\dfrac{\sqrt3}{2}\)}
    {\(\dfrac{\sqrt2}{2}\)}
    {\(\dfrac{\sqrt3}{3}\)}
    {\(\dfrac13\)}
\end{problem}

\begin{answer}
A
\end{answer}

\begin{solution}
取正方体棱长为 \(1\)，建立坐标系，使 \(A=(0,0,0)\)，\(B=(1,0,0)\)，\(C=(1,1,0)\)，\(D=(0,1,0)\)，\(A_1=(0,0,1)\)，\(B_1=(1,0,1)\)，\(D_1=(0,1,1)\). 平面 \(CB_1D_1\) 的方程为 \(x+y+z=2\)，所以过 \(A\) 且与它平行的平面 \(\alpha\) 为 \(x+y+z=0\).

与底面 \(z=0\) 相交得直线 \(m:x+y=0\)，可取方向向量 \(\bs u=(1,-1,0)\)；与侧面 \(y=0\) 相交得直线 \(n:x+z=0\)，可取方向向量 \(\bs v=(1,0,-1)\). 因此 \(\cos\angle(m,n)=\dfrac{|\bs u\cdot\bs v|}{|\bs u||\bs v|}=\dfrac{1}{2}\)，从而 \(\sin\angle(m,n)=\sqrt{1-\dfrac14}=\dfrac{\sqrt3}{2}\)，故选 A.
\end{solution}

\begin{problem}
若函数 \(f(x)=x-\dfrac13\sin 2x+a\sin x\) 在 \((-\infty,+\infty)\) 单调递增，则 \(a\) 的取值范围是（\quad）

\choices
    {\([ -1,1]\)}
    {\([ -1,\dfrac13]\)}
    {\([ -\dfrac13,\dfrac13]\)}
    {\([ -1,-\dfrac13]\)}
\end{problem}

\begin{answer}
C
\end{answer}

\begin{solution}
求导得 \(f'(x)=1-\dfrac23\cos2x+a\cos x\). 令 \(t=\cos x\)，则 \(-1\le t\le1\)，并且 \(\cos2x=2t^2-1\)，所以 \(f'(x)=\dfrac13(5+3at-4t^2)\). 关于 \(t\) 的二次函数 \(5+3at-4t^2\) 开口向下，在区间 \([-1,1]\) 上的最小值出现在端点处. 因此 \(f'(x)\ge0\) 对一切 \(x\) 成立，当且仅当 \(5+3a-4=3(a+\dfrac13)\ge0\) 且 \(5-3a-4=3(\dfrac13-a)\ge0\)，即 \(-\dfrac13\le a\le\dfrac13\)，故选 C.
\end{solution}

\section{填空题}

\begin{problem}
设向量 \(\bs a=(x,x+1)\)，\(\bs b=(1,2)\)，且 \(\bs a\perp\bs b\)，则 \(x=\fillinblank{}\).
\end{problem}

\begin{answer}
\(-\dfrac23\)
\end{answer}

\begin{solution}
向量垂直等价于数量积为零，所以 \(\bs a\cdot\bs b=x+2(x+1)=0\). 解得 \(3x+2=0\)，故 \(x=-\dfrac23\).
\end{solution}

\begin{problem}
已知 \(\theta\) 是第四象限角，且 \(\sin\!\left(\theta+\dfrac{\pi}{4}\right)=\dfrac35\)，则 \(\tan\!\left(\theta-\dfrac{\pi}{4}\right)=\fillinblank{}\).
\end{problem}

\begin{answer}
\(-\dfrac43\)
\end{answer}

\begin{solution}
令 \(\varphi=\theta+\dfrac\pi4\). 因 \(\theta\) 是第四象限角，且 \(\sin\varphi>0\)，可知 \(\varphi\) 位于第一象限，所以 \(\cos\varphi=\dfrac45\). 于是 \(\tan\varphi=\dfrac{\sin\varphi}{\cos\varphi}=\dfrac34\)，从而 \(\tan\!\left(\theta-\dfrac\pi4\right)=\tan\!\left(\varphi-\dfrac\pi2\right)=-\cot\varphi=-\dfrac43\).
\end{solution}

\begin{problem}
设直线 \(y=x+2a\) 与圆 \(C:x^2+y^2-2ay-2=0\) 相交于 \(A\)，\(B\) 两点，若 \(|AB|=2\sqrt3\)，则圆 \(C\) 的面积为\fillinblank{}.
\end{problem}

\begin{answer}
\(4\pi\)
\end{answer}

\begin{solution}
将圆方程化为 \(x^2+(y-a)^2=a^2+2\)，故圆心为 \(O=(0,a)\)，半径满足 \(r^2=a^2+2\). 直线写成 \(x-y+2a=0\)，圆心到直线的距离为 \(d=\dfrac{|a|}{\sqrt2}\). 弦长公式给出 \(|AB|^2=4(r^2-d^2)\)，所以 \(12=4\left(a^2+2-\dfrac{a^2}{2}\right)\)，解得 \(a^2=2\). 于是 \(r^2=4\)，圆 \(C\) 的面积为 \(\pi r^2=4\pi\).
\end{solution}

\begin{problem}
某高科技企业生产产品 \(A\) 和产品 \(B\)，需要甲，乙两种新型材料. 生产一件产品 \(A\) 需要甲材料 \(1.5\text{ kg}\)，乙材料 \(1\text{ kg}\)，用 \(5\) 个工时；生产一件产品 \(B\) 需要甲材料 \(0.5\text{ kg}\)，乙材料 \(0.3\text{ kg}\)，用 \(3\) 个工时，生产一件产品 \(A\) 的利润为 \(2100\) 元，生产一件产品 \(B\) 的利润为 \(900\) 元. 该企业现有甲材料 \(150\text{ kg}\)，乙材料 \(90\text{ kg}\)，则在不超过 \(600\) 个工时的条件下，生产产品 \(A\)，产品 \(B\) 的利润之和的最大值为\fillinblank{} 元.
\end{problem}

\begin{answer}
\(216000\)
\end{answer}

\begin{solution}
设生产产品 \(A\)，\(B\) 分别为 \(u\)，\(v\) 件，则约束条件为 \(\dfrac32u+\dfrac12v\le150\)，\(u+\dfrac3{10}v\le90\)，\(5u+3v\le600\)，且 \(u\ge0\)，\(v\ge0\). 化简为
\[
\begin{cases}
3u+v\le300,\\
10u+3v\le900,\\
5u+3v\le600,\\
u\ge0,\quad v\ge0.
\end{cases}
\]
利润函数为 \(P=2100u+900v\). 在 \(u\) 轴上，约束给出 \(0\le u\le90\)；在 \(v\) 轴上，约束给出 \(0\le v\le200\). 联立 \(10u+3v=900\) 与 \(5u+3v=600\)，得 \((u,v)=(60,100)\)，且满足其余约束. 另外两条约束的交点分别为 \((0,300)\) 和 \((75,75)\)，均不满足全部约束，因此可行域的顶点为 \((0,0)\)，\((90,0)\)，\((60,100)\)，\((0,200)\). 相应利润分别为 \(0\)，\(189000\)，\(216000\)，\(180000\)，所以最大利润为 \(216000\) 元.
\end{solution}

\section{解答题}

\begin{problem}
已知 \(\{a_n\}\) 是公差为 \(3\) 的等差数列，数列 \(\{b_n\}\) 满足 \(b_1=1\)，\(b_2=\dfrac13\)，\(a_nb_{n+1}+b_{n+1}=nb_n\).

\begin{enumerate}
\item 求 \(\{a_n\}\) 的通项公式；
\item 求 \(\{b_n\}\) 的前 \(n\) 项和.
\end{enumerate}
\end{problem}

\begin{solution}
\begin{enumerate}
\item 令 \(n=1\) 代入递推关系，得 \((a_1+1)b_2=b_1\)，即 \(\dfrac{a_1+1}{3}=1\)，所以 \(a_1=2\). 又数列 \(\{a_n\}\) 的公差为 \(3\)，故 \(a_n=a_1+3(n-1)=3n-1\).
\item 由 \(a_n=3n-1\) 得 \(a_n+1=3n\)，原递推关系化为 \(3n b_{n+1}=n b_n\)，所以 \(b_{n+1}=\dfrac13b_n\). 结合 \(b_1=1\)，得 \(b_n=\dfrac{1}{3^{n-1}}\). 因此 \(\{b_n\}\) 是首项为 \(1\)，公比为 \(\dfrac13\) 的等比数列，其前 \(n\) 项和为 \(S_n=\dfrac{1-(\frac13)^n}{1-\frac13}=\dfrac32-\dfrac{1}{2\cdot3^{n-1}}\).
\end{enumerate}
\end{solution}

\begin{problem}
如图，在已知正三棱锥 \(P-ABC\) 的侧面是直角三角形，\(PA=6\)，顶点 \(P\) 在平面 \(ABC\) 内的正投影为点 \(D\)，\(D\) 在平面 \(PAB\) 内的正投影为点 \(E\)，连接 \(PE\) 并延长交 \(AB\) 于点 \(G\).

\begin{enumerate}
\item 证明：\(G\) 是 \(AB\) 的中点；
\item 在图中作出点 \(E\) 在平面 \(PAC\) 内的正投影 \(F\)（说明作法及理由），并求四面体 \(PDEF\) 的体积.
\end{enumerate}

\problemresetlayout
\bitmapfigure[float=inline,max width=0.48\linewidth]{img/2016/national_paper_1_liberal/q18_fig1.png}
\end{problem}

\begin{solution}
取底面所在平面为 \(z=0\)，建立直角坐标系，令
\(A=(2\sqrt6,0,0)\)，\(B=(-\sqrt6,3\sqrt2,0)\)，\(C=(-\sqrt6,-3\sqrt2,0)\)，\(D=(0,0,0)\). 此时底面是边长为 \(6\sqrt2\) 的正三角形，且 \(D\) 是其中心. 由 \(PA=6\) 和 \(AD=2\sqrt6\)，得 \(PD=\sqrt{PA^2-AD^2}=2\sqrt3\)，取 \(P=(0,0,2\sqrt3)\).

\begin{enumerate}
\item 平面 \(PAB\) 的方程为 \(\sqrt3x+3y+\sqrt6z=6\sqrt2\)，其法向量为 \((\sqrt3,3,\sqrt6)\). 从原点 \(D\) 向该平面作垂线，垂足为
\[
E=\frac{6\sqrt2}{3+9+6}(\sqrt3,3,\sqrt6)=\left(\frac{\sqrt6}{3},\sqrt2,\frac{2\sqrt3}{3}\right)
\]
直线 \(PE\) 可写为 \(P+\lambda(E-P)\). 令其 \(z\) 坐标为零，得 \(\lambda=\dfrac32\)，所以交点
\[
G=\left(\frac{\sqrt6}{2},\frac{3\sqrt2}{2},0\right)=\frac{A+B}{2}
\]
因此 \(G\) 是 \(AB\) 的中点.
\item 在平面 \(PAB\) 内过 \(E\) 作 \(EF\parallel PB\)，并令该直线交 \(PA\) 于 \(F\). 正三棱锥给出 \(PA=PB=PC\)，而各侧面又是直角三角形，所以每个侧面的直角顶点只能是 \(P\)，从而 \(PB\perp PA\)，\(PB\perp PC\)，因此 \(PB\perp\) 平面 \(PAC\). 因此 \(EF\perp\) 平面 \(PAC\)，且 \(F\in PA\subset\) 平面 \(PAC\)，所以 \(F\) 正是点 \(E\) 在平面 \(PAC\) 内的正投影.

由 \(\overrightarrow{PB}=(-\sqrt6,3\sqrt2,-2\sqrt3)\)，可得 \(F=E-\dfrac13\overrightarrow{PB}=\left(\dfrac{2\sqrt6}{3},0,\dfrac{4\sqrt3}{3}\right)\). 于是
\[
V_{PDEF}=\frac16\left|\det\begin{pmatrix}
0&\dfrac{\sqrt6}{3}&\dfrac{2\sqrt6}{3}\\
0&\sqrt2&0\\
-2\sqrt3&-\dfrac{4\sqrt3}{3}&-\dfrac{2\sqrt3}{3}
\end{pmatrix}\right|=\frac{8}{6}=\frac43
\]
\end{enumerate}
\end{solution}

\begin{problem}
某公司计划购买 \(1\) 台机器，该种机器使用三年后即被淘汰. 机器有一易损零件，在购进机器时，可以额外购买这种零件作为备件，每个 \(200\) 元. 在机器使用期间，如果备件不足再购买，则每个 \(500\) 元. 现需决策在购买机器时应同时购买几个易损零件，为此搜集并整理了 \(100\) 台这种机器在三年试用期内更换的易损零件数，得下面柱状图：

记 \(x\) 表示 \(1\) 台机器在三年试用期内需更换的易损零件数，\(y\) 表示 \(1\) 台机器在购买易损零件上所需的费用（单位：元），\(n\) 表示购机的同时购买的易损零件数.

\begin{enumerate}
\item 若 \(n=19\)，求 \(y\) 与 \(x\) 的函数解析式；
\item 若要求“需更换的易损零件数不大于 \(n\)”的频率不小于 \(0.5\)，求 \(n\) 的最小值；
\item 假设这 \(100\) 台机器在购机的同时每台都购买 \(19\) 个易损零件，或每台都购买 \(20\) 个易损零件，分别计算这 \(100\) 台机器在购买易损零件上所需费用的平均数，以此作为决策依据，购买 \(1\) 台机器的同时应购买 \(19\) 个还是 \(20\) 个易损零件？
\end{enumerate}

\problemresetlayout
\bitmapfigure[float=inline,max width=0.70\linewidth]{img/2016/national_paper_1_liberal/q19_fig2.png}
\end{problem}

\begin{solution}
由柱状图可读出 \(x=16\)，\(17\)，\(18\)，\(19\)，\(20\)，\(21\) 时的频数分别为 \(6\)，\(16\)，\(24\)，\(24\)，\(20\)，\(10\).

\begin{enumerate}
\item 当 \(n=19\) 时，先购买的备件费用为 \(19\cdot200=3800\) 元. 当 \(x\le19\) 时备件足够，\(y=3800\)；当 \(x>19\) 时还需购买 \(x-19\) 个，每个费用为 \(500\) 元，所以
\[
y=\begin{cases}
3800,&x\le19,\\
3800+500(x-19)=500x-5700,&x>19,
\end{cases}\qquad x\in\{16,17,18,19,20,21\}
\]
\item 由频数表，需更换零件数不大于 \(18\) 的频率为 \(\dfrac{6+16+24}{100}=0.46<0.5\)，不大于 \(19\) 的频率为 \(\dfrac{6+16+24+24}{100}=0.70\ge0.5\). 所以 \(n\) 的最小值为 \(19\).
\item 同时购买 \(19\) 个时，\(x\le19\) 的 \(70\) 台每台费用为 \(3800\) 元，\(x=20\) 的 \(20\) 台每台费用为 \(4300\) 元，\(x=21\) 的 \(10\) 台每台费用为 \(4800\) 元，平均费用为 \(\dfrac{70\cdot3800+20\cdot4300+10\cdot4800}{100}=4000\) 元.

同时购买 \(20\) 个时，\(x\le20\) 的 \(90\) 台每台费用为 \(4000\) 元，\(x=21\) 的 \(10\) 台还需购买一个备件，每台费用为 \(4000+500=4500\) 元，平均费用为 \(\dfrac{90\cdot4000+10\cdot4500}{100}=4050\) 元. 因为 \(4000<4050\)，所以应同时购买 \(19\) 个易损零件.
\end{enumerate}
\end{solution}

\begin{problem}
在直角坐标系 \(xOy\) 中，直线 \(l:y=t\)（\(t\ne 0\)） 交 \(y\) 轴于点 \(M\)，交抛物线 \(C:y^2=2px\)（\(p>0\)） 于点 \(P\)，\(M\) 关于点 \(P\) 的对称点为 \(N\)，连接 \(ON\) 并延长交 \(C\) 于点 \(H\).

\begin{enumerate}
\item 求 \(\dfrac{|OH|}{|ON|}\)；
\item 除 \(H\) 以外，直线 \(MH\) 与 \(C\) 是否有其它公共点？说明理由.
\end{enumerate}
\end{problem}

\begin{solution}
\begin{enumerate}
\item 由 \(P\in C\) 且 \(y_P=t\)，得 \(P=\left(\dfrac{t^2}{2p},t\right)\)，而 \(M=(0,t)\). 由 \(P\) 是 \(MN\) 的中点，得 \(N=2P-M=\left(\dfrac{t^2}{p},t\right)\). 直线 \(ON\) 上的点可表示为 \(\lambda N\). 代入抛物线方程得 \((\lambda t)^2=2p\cdot\lambda\dfrac{t^2}{p}\)，即 \(\lambda(\lambda-2)t^2=0\). 因 \(t\ne0\)，除原点外的交点对应 \(\lambda=2\)，故 \(H=2N\)，从而 \(\dfrac{|OH|}{|ON|}=2\).
\item 由 \(H=\left(\dfrac{2t^2}{p},2t\right)\)，直线 \(MH\) 的斜率为 \(\dfrac{t}{2t^2/p}=\dfrac{p}{2t}\)，方程为 \(y=t+\dfrac{p}{2t}x\). 与 \(y^2=2px\) 联立并乘以 \(4t^2\)，得
\(\left(2t^2+px\right)^2=8pt^2x\)，\((px-2t^2)^2=0\).
因此只有一个交点 \(x=\dfrac{2t^2}{p}\)，对应 \(y=2t\)，即交点 \(H\) 为重合交点. 除 \(H\) 以外没有其它公共点.
\end{enumerate}
\end{solution}

\begin{problem}
已知函数 \(f(x)=(x-2)\e^x+a(x-1)^2\).

\begin{enumerate}
\item 讨论 \(f(x)\) 的单调性；
\item 若 \(f(x)\) 有两个零点，求 \(a\) 的取值范围.
\end{enumerate}
\end{problem}

\begin{solution}
\begin{enumerate}
\item 求导得 \(f'(x)=(x-1)\e^x+2a(x-1)=(x-1)(\e^x+2a)\).

当 \(a\ge0\) 时，\(\e^x+2a>0\)，所以 \(f'(x)<0\)（\(x<1\)），\(f'(x)>0\)（\(x>1\)），故 \(f(x)\) 在 \((-\infty,1)\) 上单调递减，在 \((1,+\infty)\) 上单调递增.

当 \(a<0\) 时，令 \(r=\ln(-2a)\)，则 \(\e^r+2a=0\). 若 \(-\dfrac\e2<a<0\)，则 \(r<1\)，于是 \(f'(x)\) 在 \((-\infty,r)\)，\((r,1)\)，\((1,+\infty)\) 上的符号依次为正、负、正，故 \(f(x)\) 在这些区间上依次单调递增、单调递减、单调递增.

若 \(a=-\dfrac\e2\)，则 \(r=1\)，且 \(f'(x)=(x-1)(\e^x-\e)\ge0\)，所以 \(f(x)\) 在 \(\R\) 上单调递增.

若 \(a<-\dfrac\e2\)，则 \(r>1\)，于是 \(f'(x)\) 在 \((-\infty,1)\)，\((1,r)\)，\((r,+\infty)\) 上的符号依次为正、负、正，故 \(f(x)\) 在这些区间上依次单调递增、单调递减、单调递增.
\item 总有 \(f(1)=-\e<0\). 当 \(a>0\) 时，\(\lim_{x\to-\infty}f(x)=+\infty\)，\(\lim_{x\to+\infty}f(x)=+\infty\)，结合第一问的单调性，\(f(x)\) 在 \((-\infty,1)\) 和 \((1,+\infty)\) 上各有且仅有一个零点.

当 \(a=0\) 时，\(f(x)=(x-2)\e^x\)，只有一个零点 \(x=2\). 当 \(a<0\) 时，\(\lim_{x\to-\infty}f(x)=-\infty\)，\(\lim_{x\to+\infty}f(x)=+\infty\). 若 \(r=\ln(-2a)\)，则
\[
f(r)=(r-2)\e^r-\frac{\e^r}{2}(r-1)^2=-\frac{\e^r}{2}\bigl((r-2)^2+1\bigr)<0
\]
在各类情形中，所有极值点的函数值均为负，所以左侧不会出现零点，而在最后一个单调递增区间上恰有一个零点. 故 \(f(x)\) 有两个零点的充要条件为 \(a>0\)，即 \(a\) 的取值范围是 \((0,+\infty)\).
\end{enumerate}
\end{solution}

\section{选考题：解答题}

\begin{problem}
如图，\(\triangle OAB\) 是等腰三角形，\(\angle AOB=120^\circ\). 以 \(O\) 为圆心，\(\dfrac12 OA\) 为半径作圆.

\begin{enumerate}
\item 证明：直线 \(AB\) 与 \(\odot O\) 相切；
\item 点 \(C\)，\(D\) 在 \(\odot O\) 上，且 \(A\)，\(B\)，\(C\)，\(D\) 四点共圆，证明：\(AB\parallel CD\).
\end{enumerate}

\bitmapfigure[max width=0.38\linewidth]{img/2016/national_paper_1_liberal/q22_fig1.png}
\end{problem}

\begin{solution}
设 \(OA=OB=R\)，以 \(O\) 为圆心、半径 \(R/2\) 的圆记为 \(\omega\).

\begin{enumerate}
\item 设 \(M\) 为 \(AB\) 的中点. 由等腰三角形的性质，\(OM\perp AB\)，且 \(\angle AOM=\dfrac12\angle AOB=60^\circ\). 在直角三角形 \(OAM\) 中，\(OM=OA\cos60^\circ=R/2\)，这正是 \(\omega\) 的半径. 因此圆心到直线 \(AB\) 的距离等于半径，故 \(AB\) 与 \(\omega\) 相切.
\item 设 \(\Gamma\) 为过 \(A\)，\(B\)，\(C\)，\(D\) 的圆，圆心为 \(Q\). 由 \(OA=OB\)，点 \(O\) 在 \(AB\) 的垂直平分线上；又因 \(QA=QB\)，点 \(Q\) 也在 \(AB\) 的垂直平分线上. 因此直线 \(OQ\perp AB\). 两圆 \(\omega\) 和 \(\Gamma\) 的公共弦为 \(CD\)，公共弦垂直于两圆心的连线，所以 \(CD\perp OQ\). 于是 \(CD\parallel AB\).
\end{enumerate}
\end{solution}

\begin{problem}

在直角坐标系 \(xOy\) 中，曲线 \(C_1\) 的参数方程为
\(\begin{cases}
x=a\cos t,\\
y=1+a\sin t,
\end{cases}\)
（\(t\) 为参数，\(a>0\)）.
在以坐标原点为极点，\(x\) 轴正半轴为极轴的极坐标系中，曲线 \(C_2:\rho=4\cos\theta\).

\begin{enumerate}
\item 说明 \(C_1\) 是哪种曲线，并将 \(C_1\) 的方程化为极坐标方程；
\item 直线 \(C_3\) 的极坐标方程为 \(\theta=\alpha_0\)，其中 \(\alpha_0\) 满足 \(\tan\alpha_0=2\)，若曲线 \(C_1\) 与 \(C_2\) 的公共点都在 \(C_3\) 上，求 \(a\).
\end{enumerate}
\end{problem}

\begin{solution}
\begin{enumerate}
\item 由参数方程得 \(x^2=(a\cos t)^2\)，\((y-1)^2=(a\sin t)^2\)，相加得 \(x^2+(y-1)^2=a^2\). 所以 \(C_1\) 是圆心为 \((0,1)\)、半径为 \(a\) 的圆. 利用 \(x=\rho\cos\theta\)，\(y=\rho\sin\theta\)，其极坐标方程为
\[
\rho^2\cos^2\theta+(\rho\sin\theta-1)^2=a^2
\]
即 \(\rho^2-2\rho\sin\theta+1-a^2=0\).
\item 取曲线 \(C_1\) 与 \(C_2\) 的一个公共点. 由题设该点在 \(C_3\) 上，故可取其极角为 \(\theta=\alpha_0\)，并且由 \(C_2\) 得 \(\rho=4\cos\alpha_0\). 代入 \(C_1\) 的极坐标方程，得
\[
16\cos^2\alpha_0-8\sin\alpha_0\cos\alpha_0+1-a^2=0
\]
由 \(\tan\alpha_0=2\)，有 \(\cos^2\alpha_0=\dfrac15\)，\(\sin\alpha_0\cos\alpha_0=\dfrac25\)，代入上式得 \(\dfrac{16}{5}-\dfrac{16}{5}+1-a^2=0\). 因 \(a>0\)，所以 \(a=1\).
\end{enumerate}
\end{solution}

\begin{problem}

已知函数 \(f(x)=|x+1|-|2x-3|\).

\begin{enumerate}
\item 在图中画出 \(y=f(x)\) 的图象；
\item 求不等式 \(|f(x)|>1\) 的解集.
\end{enumerate}

\bitmapfigure[max width=0.45\linewidth]{img/2016/national_paper_1_liberal/q24_fig1.png}
\end{problem}

\begin{solution}
\begin{enumerate}
\item 按照两个绝对值内式子的零点 \(-1\) 和 \(\dfrac32\) 分段，得
\[
f(x)=
\begin{cases}
x-4,&x\le-1,\\
3x-2,&-1<x\le\dfrac32,\\
-x+4,&x>\dfrac32.
\end{cases}
\]
所以图象由三段直线组成：\(x\le-1\) 时为直线 \(y=x-4\)，\(-1<x\le\dfrac32\) 时为直线 \(y=3x-2\)，\(x>\dfrac32\) 时为直线 \(y=-x+4\)，并在 \((-1,-5)\) 和 \(\left(\dfrac32,\dfrac52\right)\) 处连接.
\item 分段解不等式. 对于 \(x\le-1\)，有 \(f(x)=x-4<-1\)，故该区间全部满足. 对于 \(-1<x\le\dfrac32\)，有 \(|3x-2|>1\)，所以 \(x<\dfrac13\) 或 \(x>1\)，与本区间取交得 \((-1,\dfrac13)\cup(1,\dfrac32]\). 对于 \(x>\dfrac32\)，有 \(|4-x|>1\)，所以 \(x<3\) 或 \(x>5\)，与本区间取交得 \((\dfrac32,3)\cup(5,+\infty)\). 合并各段，解集为 \(\left\{x\mid x<\dfrac13\text{ 或 }1<x<3\text{ 或 }x>5\right\}\).
\end{enumerate}
\end{solution}
